The altitude to the length-$10$ base is $12$, so the area is $60$ and the inradius is $r=\dfrac{60}{18}=\dfrac{10}{3}$. The apex corner triangle has altitude $12-2r=\dfrac{16}{3}$, hence similarity ratio $\dfrac{4}{9}$ and area $\dfrac{320}{27}$. The altitude to either equal side is $\dfrac{120}{13}$. Each base-corner triangle therefore has similarity ratio $\dfrac{\dfrac{120}{13}-\dfrac{20}{3}}{\dfrac{120}{13}}=\dfrac{5}{18}$ and area $\dfrac{125}{27}$. The hexagon area is $60-\dfrac{320+125+125}{27}=\dfrac{350}{9}=38+\dfrac{8}{9}$. Thus $\boxed{m+n+l=55}$.
\begin{asy}
texpreamble("\usepackage{amsmath,amssymb}");
size(240);
defaultpen(fontsize(9pt));
pair A=(0,12),B=(-5,0),C=(5,0);draw(A--B--C--cycle);filldraw((-20/9.0,20/3.0)--(20/9.0,20/3.0)--(65/18.0,10/3.0)--(20/9.0,0)--(-20/9.0,0)--(-65/18.0,10/3.0)--cycle,lightblue);draw(circle((0,10/3.0),10/3.0));label("A",A,N);label("B",B,SW);label("C",C,SE);label("13",(A+B)/2,W);label("13",(A+C)/2,E);label("10",(B+C)/2,S);
\end{asy}
